\section{Conclusion}
\begingroup
Neural eigenfluids express a physically structured fluid spectrum through a continuous spatial basis. Differential lifts and potential traces define admissible velocities; geometry-local orthogonalization and Rayleigh minimization organize their positive modes. Canonical gauge selection gives potential-space learning a lower-derivative formulation, with the spectral-variance relation explaining its weighting of approximate modes. A neural Dirichlet construction supplies the independent harmonic circulation coordinates.

The experiments connect this structure to representation and dynamics. Learned spectral spans sustain accurate periodic vortex merger and improve energy capture and autonomous accuracy over equal-capacity Laplacian controls on directly sampled random fields. Completed bases express obstacle-scale circulation, and matched conditioning tests identify how harmonic coordinates affect evolution. In three dimensions, increasing capacity reduces spatial representation error while exposing the contribution of coefficient dynamics to total flow error. Full quadratic assembly moves the same reduced generator into coefficient space, allowing repeated initial states to share their spatial basis and operators. Continuous evaluation, physical spectral organization, and reusable autonomous computation are thus realized within one reduced representation.
\par\endgroup
